\markdownRendererDocumentBegin
\markdownRendererUlBeginTight
\markdownRendererUlItem change point analysis for each region?\markdownRendererInterblockSeparator
{}\markdownRendererUlBeginTight
\markdownRendererUlItem does a stat sig. chng pt negate a stat sig trend?\markdownRendererUlItemEnd 
\markdownRendererUlEndTight \markdownRendererUlItemEnd 
\markdownRendererUlItem should we subset recent data to match frequency of older data? or should we repeate analysis using the daily interpolated data?\markdownRendererUlItemEnd 
\markdownRendererUlItem include model data (or mooring data) for subregions stratification or heat content? this is a biiig can of worms\markdownRendererUlItemEnd 
\markdownRendererUlItem mooring data might be doable but limited regions and years\markdownRendererUlItemEnd 
\markdownRendererUlItem calculate first last and ndays using the ACTUAL eqn you show\markdownRendererUlItemEnd 
\markdownRendererUlItem also calculate the DAY OF MAX\markdownRendererUlItemEnd 
\markdownRendererUlItem is some form of trend analysis possible using a subset of the data?\markdownRendererInterblockSeparator
{}\markdownRendererUlBeginTight
\markdownRendererUlItem specifically, if we categorize high/med/low ice years can we perform linear regresssion on each set?\markdownRendererUlItemEnd 
\markdownRendererUlItem seems like discontinuity (e.g. varying \markdownRendererDollarSign{}\markdownRendererBackslash{}Delta t\markdownRendererDollarSign{}) would be a problem\markdownRendererUlItemEnd 
\markdownRendererUlEndTight \markdownRendererUlItemEnd 
\markdownRendererUlEndTight \markdownRendererDocumentEnd